\begingroup
\renewcommand{\arraystretch}{1.45}
\SourceLeaf{302}
\noindent
\(m^{\mathrm{tum}}\) deprimatur, unde necessario functio \(X\) ipsa, quae e valoribus \(A,B,C,D\) etc. 
per methodum art. praec. eruitur, prodire debebit. Quum methodus terminum 
summum functionis \(Y\) atque inde deinceps inferiores evolvendi, tamquam casus 
specialis e methodo in art. 15, II tradita peti possit, non opus est, huic rei hic 
diutius inhaerere.

\InterpolationArticle{18}
Omnes disquisitiones praecedentes superstructae sunt conditioni, \(X\) esse progressionem cum cosinu sinuque arcus \(mx\) abruptam. Quae conditio si locum non 
habet, tot valores ipsius \(X\), quot cognitos esse hactenus assumsimus, manifesto non 
sufficiunt ad determinationem completam, expressioque pro \(T\) in variis quos consideravimus casibus tradita, correctione opus habebit, a co\"efficientibus sequentibus 
in \(X\) pendente et per summationes artt. 1 et 2 facile assignabili. Quodsi autem 
series \(X\), sive in infinitum excurrat sive uspiam abrumpatur, tantopere convergit, 
ut partes post \(\cos mx\) et \(\sin mx\) sequentes pro evanescentibus haberi possint, 
etiam illam correctionem negligere licebit, adeoque valorem \(T\) saltem proxime 
verum nacti erimus. Quomodocunque vero haec se habeant, \(T\) certo functionem 
verae similem simplicissimam exhibet, per quam omnibus valoribus propositis satisfieri potest.

\InterpolationArticle{19}
Imprimis frequens, maximaque adeo attentione dignus est casus iste, ubi 
valores arcus \(x\), pro quibus valores functionis \(X\) dati sunt, progressionem arithmeticam constituunt, in qua terminorum differentia aequalis est co\"efficienti, ex 
divisione peripheriae integrae in totidem partes quot sunt termini orto, ita ut si 
progressio ultra terminum postremum continuaretur, termini primi peripheria integra aucti rursus deinceps prodituri essent. Postquam itaque disquisitionis generalis praecipua momenta in praecedentibus absolvimus, huncce casum, ubi 
termini dati quasi \emph{periodum} completam constituunt, coronidis loco seorsim copiosius tractabimus. Antequam vero hanc disquisitionem ipsam aggrediamur, duo 
lemmata erunt praemittenda.

\textsc{Lemma Primum.} \emph{Si arcus \(a,b,c,d\) etc., quorum multitudo est \(\mu\), constituunt 
progressionem arithmeticam, in qua terminorum differentia \(b-a=c-b=d-c\) 
etc. est \(=\frac{360^{0}}{\mu}\), productum ex \(\mu\) factoribus \(P=\)}
\SourceLeaf{303}
\[
\begin{aligned}
&\sin\frac12(t-a)\sin\frac12(t-b)\sin\frac12(t-c)\sin\frac12(t-d)\ \mathrm{etc.}
\end{aligned}
\]
\noindent\emph{fit \(=\mp\frac{\sin\frac12\mu(t-a)}{2^{\mu-1}}\), ubi signum superius vel inferius valet, prout \(\mu\) par est vel impar.}

\emph{Demonstr.} Potest quidem hoc theorema facile ex alio ab ill. \textsc{Eulero} in \emph{Introd. in Anal. Infin.} I, \S. 240 tradito derivari; ne quid vero hic desiderari videatur, demonstrationem aliam paucis hic explicabimus.

I. Si productum \(P\) primo ponitur sub formam
\[
\begin{aligned}
&\cos\frac12(t-a-180^{0})\cos\frac12(t-b-180^{0})\\[5pt]&\qquad\cos\frac12(t-c-180^{0})\cos\frac12(t-d-180^{0})\mathinner{.\,.\,.}
\end{aligned}
\]
\noindent atque tunc secundum algorithmum sinuum notum in aggregatum cosinuum evolvitur, facile perspicietur, inde prodire functionem vel formae \(F\) vel formae \(G\) 
(art. 13) ad ordinem \(\frac12\mu^{\mathrm{tum}}\) ascendentem, terminumque \emph{summum} fieri
\[
\begin{aligned}
&=\frac{1}{2^{\mu-1}}\cos\frac12(\mu t-a-b-c-d-\mathrm{etc.}-\mu\times180^{0})\\[5pt]
&=\frac{1}{2^{\mu-1}}\cos\bigl(\frac12\mu(t-a-360^{0})+90^{0}\bigr)\\[5pt]
&=-\frac{1}{2^{\mu-1}}\sin\frac12\mu(t-a-360^{0})\\[5pt]
&=-\frac{1}{2^{\mu-1}}\sin\frac12\mu(t-a)\cos\mu\times180^{0}\\[5pt]
&=\mp\frac{1}{2^{\mu-1}}\sin\frac12\mu(t-a)
\end{aligned}
\]
\noindent signo superiore pro valore pari ipsius \(\mu\) valente, inferiore pro impari.

II. Quum \(\sin\frac12\mu(t-a)\) fiat \(=0\), si ipsi \(t\) aliquis valorum \(a.\ b,c,d\) etc. 
tribuitur, per theorema art. 14 erit \(\sin\frac12\mu(t-a)\) per productum \(P\) divisibilis; 
quotiens autem necessario erit quantitas definita \(=h\), quoniam \(\sin\frac12\mu(t-a)\) et \(P\) 
eiusdem sunt ordinis. Manifesto autem (per I) ex aequatione \(P=\frac1h\sin\frac12\mu(t-a)\) 
sequitur esse \(h=\mp2^{\mu-1}\), adeoque \(P=\mp\frac{\sin\frac12\mu(t-a)}{2^{\mu-1}}\)\quad Q.~E.~D.

\textsc{Lemma Secundum.} \emph{Si arcus \(a,b,c,d\mathinner{.\,.\,.}\) eodem modo se habent, ut in lemmate primo, productum ex \(\mu\) factoribus}
\[
\begin{aligned}
&(\cos t-\cos a)(\cos t-\cos b)(\cos t-\cos c)(\cos t-\cos d)\mathinner{.\,.\,.}
\end{aligned}
\]
\noindent\emph{fit \(=\frac{\cos\mu t-\cos\mu a}{2^{\mu-1}}\)}

\emph{Demonstr.} Quum fiat
\SourceLeaf{304}
\[
\begin{aligned}
&\cos t-\cos a=2\sin\frac12(t-a)\sin\frac12(-t-a)\\[5pt]
&\cos t-\cos b=2\sin\frac12(t-b)\sin\frac12(-t-b)\\[5pt]
&\cos t-\cos c=2\sin\frac12(t-c)\sin\frac12(-t-c)\\[5pt]
&\cos t-\cos d=2\sin\frac12(t-d)\sin\frac12(-t-d)\\[5pt]
&\mathrm{etc.}
\end{aligned}
\]
\noindent productum ex his factoribus fit, per lemma primum
\[
\begin{aligned}
&=2^{\mu}\times\mp\frac{\sin\frac12\mu(t-a)}{2^{\mu-1}}\times\mp\frac{\sin\frac12\mu(-t-a)}{2^{\mu-1}}\\[5pt]
&=\frac{\sin\frac12\mu(t-a)\sin\frac12\mu(-t-a)}{2^{\mu-2}}=\frac{\cos\mu t-\cos\mu a}{2^{\mu-1}}\quad\mathrm{Q.\ E.\ D.}
\end{aligned}
\]

Quum habeatur
\[
\begin{aligned}
&\sin\frac12\mu(t-a)=-\sin\frac12\mu(t-b)\\[5pt]&=\sin\frac12\mu(t-c)=-\sin\frac12\mu(t-d)\ \mathrm{etc.}
\end{aligned}
\]
\noindent nec non \(\cos\mu a=\cos\mu b=\cos\mu c=\cos\mu d\) etc.: manifestum est, productum 
in lemmate primo fieri etiam
\[
\begin{aligned}
&=\pm\frac{\sin\frac12\mu(t-b)}{2^{\mu-1}}\\[5pt]&=\mp\frac{\sin\frac12\mu(t-c)}{2^{\mu-1}}=\pm\frac{\sin\frac12\mu(t-d)}{2^{\mu-1}}\ \mathrm{etc.}
\end{aligned}
\]
\noindent nec non productum in lemmate secundo fieri etiam
\[
\begin{aligned}
&=\frac{\cos\mu t-\cos\mu b}{2^{\mu-1}}=\frac{\cos\mu t-\cos\mu c}{2^{\mu-1}}\\[5pt]&=\frac{\cos\mu t-\cos\mu d}{2^{\mu-1}}\ \mathrm{etc.}
\end{aligned}
\]

\InterpolationArticle{20}
Consideremus primo casum generalem art. 10, ubi \(X\) est formae
\[
\begin{aligned}
&\alpha+\alpha'\cos x+\alpha''\cos2x+\alpha'''\cos3x+\mathinner{.\,.\,.}+\alpha^{(m)}\cos mx\\[5pt]
&\quad+\beta'\sin x+\beta''\sin2x+\beta'''\sin3x+\mathinner{.\,.\,.}+\beta^{(m)}\sin mx
\end{aligned}
\]
\noindent atque \(\mu=2m+1\). Hic igitur erit
\[
\begin{aligned}
&\sin\frac12(t-b)\sin\frac12(t-c)\sin\frac12(t-d)\mathinner{.\,.\,.}=\frac{\sin\frac12\mu(t-a)}{2^{\mu-1}\sin\frac12(t-a)}\\[5pt]
&=\frac{1}{2^{\mu-1}}\bigl(1+2\cos(t-a)+2\cos2(t-a)\\[5pt]&\qquad{}+2\cos3(t-a)+\mathinner{.\,.\,.}+2\cos m(t-a)\bigr)
\end{aligned}
\]
\noindent unde substituendo \(a\) pro \(t\)
\[
\begin{aligned}
&\sin\frac12(a-b)\sin\frac12(a-c)\sin\frac12(a-d)\mathinner{.\,.\,.}=\frac{\mu}{2^{\mu-1}}
\end{aligned}
\]
\SourceLeaf{305}
Hinc in formula art. 10 pro \(T\), fit co\"efficiens ipsius \(A\)
\[
\begin{aligned}
&=\frac{1}{\mu}(1+2\cos(t-a)+2\cos2(t-a)\\[5pt]&\qquad{}+2\cos3(t-a)+\mathinner{.\,.\,.}+2\cos m(t-a))
\end{aligned}
\]
Simili modo fit
\[
\begin{aligned}
&\sin\frac12(t-a)\sin\frac12(t-c)\sin\frac12(t-d)\mathinner{.\,.\,.}
=-\frac{\sin\frac12\mu(t-b)}{2^{\mu-1}\sin\frac12(t-b)}\\[5pt]
&=-\frac{1}{2^{\mu-1}}(1+2\cos(t-b)+2\cos2(t-b)\\[5pt]&\qquad{}+2\cos3(t-b)+\mathinner{.\,.\,.}+2\cos m(t-b))
\end{aligned}
\]
atque
\[
\begin{aligned}
&\sin\frac12(b-a)\sin\frac12(b-c)\sin\frac12(b-d)\mathinner{.\,.\,.}=-\frac{\mu}{2^{\mu-1}}
\end{aligned}
\]
Hinc co\"efficiens ipsius \(B\) in formula pro \(T\) fit
\[
\begin{aligned}
&=\frac{1}{\mu}(1+2\cos(t-b)+2\cos2(t-b)\\[5pt]&\qquad{}+2\cos3(t-b)+\mathinner{.\,.\,.}+2\cos m(t-b))
\end{aligned}
\]
Prorsus similes erunt co\"efficientes ipsorum \(C,D\) etc., unde tandem concluditur,
fieri
\[
\begin{aligned}
&T=\frac{1}{\mu}(A+B+C+D+\mathinner{.\,.\,.})\\[5pt]
&+\frac{2}{\mu}(A\cos a+B\cos b+C\cos c+D\cos d+\mathinner{.\,.\,.})\cos t\\[5pt]
&+\frac{2}{\mu}(A\sin a+B\sin b+C\sin c+D\sin d+\mathinner{.\,.\,.})\sin t\\[5pt]
&+\frac{2}{\mu}(A\cos2a+B\cos2b+C\cos2c+D\cos2d+\mathinner{.\,.\,.})\cos2t\\[5pt]
&+\frac{2}{\mu}(A\sin2a+B\sin2b+C\sin2c+D\sin2d+\mathinner{.\,.\,.})\sin2t\\[5pt]
&+\mathrm{etc.}\\[5pt]
&+\frac{2}{\mu}(A\cos ma+B\cos mb+C\cos mc+D\cos md+\mathinner{.\,.\,.})\cos mt\\[5pt]
&+\frac{2}{\mu}(A\sin ma+B\sin mb+C\sin mc+D\sin md+\mathinner{.\,.\,.})\sin mt
\end{aligned}
\]
Quum haec formula cum
\[
\begin{aligned}
&\alpha+\alpha'\cos t+\beta'\sin t+\alpha''\cos2t+\beta''\sin2t+\mathrm{etc.}\\[5pt]&\qquad{}+\alpha^{(m)}\cos mt+\beta^{(m)}\sin mt
\end{aligned}
\]
identica esse debeat, valores co\"efficientium \(\alpha,\alpha',\beta',\alpha'',\beta''\) etc. hinc protinus
habentur.
\InterpolationArticle{21}
Si progressio pro \(X\) cum terminis \(\cos mx\) et \(\sin mx\) non abrumpitur, sive
in infinitum excurrat, sive finita sit, valor pro \(T\) in art. praec. inventus incom\SourceLeaf{306}%
pletus erit, itemque valores singulorum co\"efficientium \(\alpha,\alpha',\beta',\alpha'',\beta''\) etc. ibi
traditi correctione opus habebunt. Haec autem commodius per methodum sequentem quam per summationem art. 2 determinatur. Ante omnia observamus, esse
\[
\begin{aligned}
&\cos na+\cos nb+\cos nc+\cos nd+\mathinner{.\,.\,.}=\mu\cos na\\[5pt]
&\sin na+\sin nb+\sin nc+\sin nd+\mathinner{.\,.\,.}=\mu\sin na
\end{aligned}
\]
quoties \(n\) est integer per \(\mu\) divisibilis; contra
\[
\begin{aligned}
&\cos na+\cos nb+\cos nc+\cos nd+\mathinner{.\,.\,.}=0\\[5pt]
&\sin na+\sin nb+\sin nc+\sin nd+\mathinner{.\,.\,.}=0
\end{aligned}
\]
quoties \(n\) est integer per \(\mu\) non divisibilis. Pro casu priore res per se clara est;
pro posteriore sit summa prima \(=P\), secunda \(=Q\), unde facile deducitur
\[
\begin{aligned}
&P\cos n(b-a)-Q\sin n(b-a)=P\\[5pt]
&P\sin n(b-a)+Q\cos n(b-a)=Q
\end{aligned}
\]
Multiplicando aequationem primam per \(\cos n(b-a)-1\), secundam per \(\sin n(b-a)\),
fit addendo delendoque quae mutuo se destruunt
\[
\begin{aligned}
&2P(1-\cos n(b-a))=0
\end{aligned}
\]
Similiter multiplicando aequationem primam per \(\sin n(b-a)\), secundam per
\(1-\cos n(b-a)\), provenit addendo
\[
\begin{aligned}
&2Q(1-\cos n(b-a))=0
\end{aligned}
\]
Iam pro casu quidem priore (ubi \(n\) per \(\mu\) divisibilis est, et proin \(\cos n(b-a)=1\))
hae aequationes identicae sunt, pro posteriore autem (ubi \(n\) per \(\mu\) non est divisibilis, adeoque \(\cos n(b-a)=\cos\frac{n}{\mu}\times360^{0}\) non potest esse \(=1\)) consistere nequeunt, nisi fuerit \(P=0\) et \(Q=0\).

Iam ponamus, post terminos \(\alpha^m\cos mx+\beta^m\sin mx\) in expressione functionis \(X\) sequi
\[
\begin{aligned}
&\phantom{{}+{}}\alpha^{m+1}\cos(m+1)x+\alpha^{m+2}\cos(m+2)x+\mathrm{etc}\\[5pt]
&+\beta^{m+1}\sin(m+1)x+\beta^{m+2}\sin(m+2)x+\mathrm{etc.}
\end{aligned}
\]
valorem vero (incompletum), in suppositione, has partes non adesse, pro \(T\) in
art. praec. inventum, esse
\SourceLeaf{307}
\[
\begin{aligned}
&\gamma+\gamma'\cos t+\gamma''\cos2t+\gamma'''\cos3t+\mathinner{.\,.\,.}+\gamma^m\cos mt\\[5pt]
&\phantom{\gamma}\;+\delta'\sin t+\delta''\sin2t+\delta'''\sin3t+\mathinner{.\,.\,.}+\delta^m\sin mt
\end{aligned}
\]
ita ut habeatur \(\gamma=\frac{1}{\mu}(A+B+C+D+\mathrm{etc.})\) etc. Substituantur pro \(A,B,C,D\)
etc. valores sui, puta
\[
\begin{aligned}
&A=\alpha+\alpha'\cos a+\alpha''\cos2a+\mathinner{.\,.\,.}+\alpha^m\cos ma\\[5pt]&\qquad{}+\alpha^{m+1}\cos(m+1)a\mathinner{.\,.\,.}\\[5pt]
&+\beta'\sin a+\beta''\sin2a+\mathinner{.\,.\,.}+\beta^m\sin ma\\[5pt]&\qquad{}+\beta^{m+1}\sin(m+1)a\mathinner{.\,.\,.}\\[5pt]
&\mathrm{etc.}
\end{aligned}
\]
fietque (per summationem modo traditam)
\[
\begin{aligned}
&\gamma=\alpha+\alpha^{\mu}\cos\mu a+\alpha^{2\mu}\cos2\mu a+\alpha^{3\mu}\cos3\mu a+\mathrm{etc.}\\[5pt]
&+\beta^{\mu}\sin\mu a+\beta^{2\mu}\sin2\mu a+\beta^{3\mu}\sin3\mu a+\mathrm{etc.}\\[5pt]
&\gamma'=\alpha'+(\alpha^{\mu-1}+\alpha^{\mu+1})\cos\mu a\\[5pt]&\qquad{}+(\alpha^{2\mu-1}+\alpha^{2\mu+1})\cos2\mu a+\mathrm{etc.}\\[5pt]
&+(\beta^{\mu-1}+\beta^{\mu+1})\sin\mu a\\[5pt]&\qquad{}+(\beta^{2\mu-1}+\beta^{2\mu+1})\sin2\mu a+\mathrm{etc.}\\[5pt]
&\delta'=\beta'-(\beta^{\mu-1}-\beta^{\mu+1})\cos\mu a\\[5pt]&\qquad{}-(\beta^{2\mu-1}-\beta^{2\mu+1})\cos2\mu a-\mathrm{etc.}\\[5pt]
&-(\alpha^{\mu-1}-\alpha^{\mu+1})\sin\mu a\\[5pt]&\qquad{}-(\alpha^{2\mu-1}-\alpha^{2\mu+1})\sin2\mu a-\mathrm{etc.}\\[5pt]
&\gamma''=\alpha''+(\alpha^{\mu-2}+\alpha^{\mu+2})\cos\mu a\\[5pt]&\qquad{}+(\alpha^{2\mu-2}+\alpha^{2\mu+2})\cos2\mu a+\mathrm{etc.}\\[5pt]
&+(\beta^{\mu-2}+\beta^{\mu+2})\sin\mu a\\[5pt]&\qquad{}+(\beta^{2\mu-2}+\beta^{2\mu+2})\sin2\mu a+\mathrm{etc.}\\[5pt]
&\delta''=\beta''-(\beta^{\mu-2}-\beta^{\mu+2})\cos\mu a\\[5pt]&\qquad{}-(\beta^{2\mu-2}-\beta^{2\mu+2})\cos2\mu a-\mathrm{etc.}\\[5pt]
&-(\alpha^{\mu-2}-\alpha^{\mu+2})\sin\mu a\\[5pt]&\qquad{}-(\alpha^{2\mu-2}-\alpha^{2\mu+2})\sin2\mu a-\mathrm{etc.}
\end{aligned}
\]
etc. usque ad
\[
\begin{aligned}
&\gamma^m=\alpha^m+(\alpha^{\mu-m}+\alpha^{\mu+m})\cos\mu a\\[5pt]&\qquad{}+(\alpha^{2\mu-m}+\alpha^{2\mu+m})\cos2\mu a+\mathrm{etc.}\\[5pt]
&+(\beta^{\mu-m}+\beta^{\mu+m})\sin\mu a\\[5pt]&\qquad{}+(\beta^{2\mu-m}+\beta^{2\mu+m})\sin2\mu a+\mathrm{etc.}\\[5pt]
&\delta^m=\beta^m-(\beta^{\mu-m}-\beta^{\mu+m})\cos\mu a\\[5pt]&\qquad{}-(\beta^{2\mu-m}-\beta^{2\mu+m})\cos2\mu a-\mathrm{etc.}\\[5pt]
&-(\alpha^{\mu-m}-\alpha^{\mu+m})\sin\mu a\\[5pt]&\qquad{}-(\alpha^{2\mu-m}-\alpha^{2\mu+m})\sin2\mu a-\mathrm{etc.}
\end{aligned}
\]
Si itaque progressio tantopere convergit. ut \(\alpha^{m+1},\beta^{m+1},\alpha^{m+2},\beta^{m+2}\) etc. negligi
possint, valores \(\gamma,\gamma',\delta'\) etc. pro veris \(\alpha,\alpha',\beta'\) etc. accipere licebit.
\InterpolationArticle{22}
Transimus ad casum alterum art. 11, ubi progressio \(X\) quidem eadem manet ut in art. 20, multitudo valorum datorum autem unitate minor est, quam
multitudo co\"efficientium incognitorum, puta \(=2m=\mu\). Hic habetur
\SourceLeaf{308}
\[
\begin{aligned}
&\cos\frac12(t-a)\sin\frac12(t-b)\sin\frac12(t-c)\sin\frac12(t-d)\mathinner{.\,.\,.}\\[5pt]
&=-\frac{\sin\frac12\mu(t-a)\cos\frac12(t-a)}{2^{\mu-1}\sin\frac12(t-a)}\\[5pt]
&=-\frac{\cos\frac12(t-a)}{2^{\mu-2}}\bigl(\cos\frac12(t-a)+\cos\frac32(t-a)\\[5pt]&\qquad{}+\cos\frac52(t-a)+\mathinner{.\,.\,.}+\cos\frac{\mu-1}{2}(t-a)\bigr)\\[5pt]
&=-\frac{1}{2^{\mu-1}}\bigl(1+2\cos(t-a)+2\cos2(t-a)\\[5pt]&\qquad{}+2\cos3(t-a)+\mathinner{.\,.\,.}+2\cos(m-1)(t-a)\\[5pt]
&\qquad +\cos m(t-a)\bigr)
\end{aligned}
\]
Quare in formula art. 11 pro $T$ co\"efficiens ipsius $A$ fit
\[
\begin{aligned}
&=\frac1\mu\bigl(1+2\cos(t-a)+2\cos2(t-a)\\[5pt]&\qquad{}+2\cos3(t-a)+\mathinner{.\,.\,.}\\[5pt]
&\qquad +2\cos(m-1)(t-a)+\cos m(t-a)\bigr)
\end{aligned}
\]
Et prorsus similes expressiones (mutato tantummodo $a$ in $b,c,d$ etc.) pro co\"efficientibus ipsorum $B,C,D$ etc. inveniuntur. Denique pars ultima ipsius $T$
\[
\begin{aligned}
&2^{2m-1}i^{2m}\sin\frac12(t-a)\sin\frac12(t-b)\sin\frac12(t-c)\sin\frac12(t-d)\mathinner{.\,.\,.}\\[5pt]&\qquad\times\{\alpha^m\cos\frac12s+\beta^m\sin\frac12s\}
\end{aligned}
\]
fit
\[
\begin{aligned}
&=-i^{2m}\sin m(t-a)\{\alpha^m\cos(ma+(2m-1)90^{0})\\[5pt]&\qquad{}+\beta^m\sin(ma+(2m-1)90^{0})\}\\[5pt]
&=-\sin m(t-a)(\alpha^m\sin ma-\beta^m\cos ma)
\end{aligned}
\]
Hinc tandem colligitur
\[
\begin{aligned}
&T=\frac1\mu(A+B+C+D+\mathrm{etc.})\\[5pt]
&\quad+\frac2\mu(A\cos a+B\cos b+C\cos c+D\cos d+\mathrm{etc.})\cos t\\[5pt]
&\quad+\frac2\mu(A\sin a+B\sin b+C\sin c+D\sin d+\mathrm{etc.})\sin t\\[5pt]
&\quad+\frac2\mu(A\cos2a+B\cos2b+C\cos2c+D\cos2d+\mathrm{etc.})\cos2t\\[5pt]
&\quad+\frac2\mu(A\sin2a+B\sin2b+C\sin2c+D\sin2d+\mathrm{etc.})\sin2t\\[5pt]
&\quad+\mathrm{etc.}\\[5pt]
&\quad+\frac1\mu(A\cos ma+B\cos mb+C\cos mc+D\cos md+\mathrm{etc.})\cos mt\\[5pt]
&\quad+\frac1\mu(A\sin ma+B\sin mb+C\sin mc+D\sin md+\mathrm{etc.})\sin mt\\[5pt]
&\quad+(\alpha^m\sin ma-\beta^m\cos ma)\\[5pt]&\qquad(\sin ma\cos mt-\cos ma\sin mt)
\end{aligned}
\]
Quae expressio, si, rescissa parte ultima, per
\[
\begin{aligned}
&\gamma+\gamma'\cos t+\gamma''\cos2t+\mathrm{etc.}+\gamma^m\cos mt\\[5pt]
&\quad+\delta'\sin t+\delta''\sin2t+\mathrm{etc.}+\delta^m\sin mt
\end{aligned}
\]
\SourceLeaf{309}
exhibetur, omnes co\"efficientes $\gamma,\gamma',\delta',\gamma'',\delta''$ etc. praecise eosdem valores habere patet, ut in art. praec., ubi erat $\mu=2m+1$, exceptis duobus ultimis $\gamma^m,\delta^m$, qui illic duplo erant maiores. Porro liquet, omnes co\"efficientes
$\alpha,\alpha',\alpha''\mathinner{.\,.\,.}\alpha^{m-1},\beta',\beta''\mathinner{.\,.\,.}\beta^{m-1}$ illis $\gamma,\gamma',\gamma''\mathinner{.\,.\,.}\gamma^{m-1},\delta',\delta''\mathinner{.\,.\,.}\delta^{m-1}$ resp. aequales adeoque cognitos fieri, inter ultimos $\alpha^m,\beta^m$ autem haberi aequationes
\[
\begin{aligned}
&\alpha^m=\gamma^m+(\alpha^m\sin ma-\beta^m\cos ma)\sin ma\\[5pt]
&\beta^m=\delta^m-(\alpha^m\sin ma-\beta^m\cos ma)\cos ma
\end{aligned}
\]
quae tamen (ut ex natura problematis praevidere licebat) non sufficiunt ad determinationem completam co\"efficientium $\alpha^m,\beta^m$. Scilicet quum fiat
\[
\begin{aligned}
&\cos ma=-\cos mb=\cos mc=-\cos md\ \mathrm{etc.}
\end{aligned}
\]
atque
\[
\begin{aligned}
&\sin ma=-\sin mb=\sin mc=-\sin md\ \mathrm{etc.}
\end{aligned}
\]
adeoque
\[
\begin{aligned}
&\gamma^m=\frac{\cos ma}{\mu}(A-B+C-D+\mathrm{etc.})\\[5pt]
&\delta^m=\frac{\sin ma}{\mu}(A-B+C-D+\mathrm{etc.})
\end{aligned}
\]
facile perspicitur, utramque aequationem contentam esse sub unica
\[
\begin{aligned}
&\alpha^m\cos ma+\beta^m\sin ma=\frac1\mu(A-B+C-D+\mathrm{etc.})
\end{aligned}
\]
Aliam itaque insuper relationem inter $\alpha^m$ et $\beta^m$ datam esse oportet, ut utrumque co\"efficientem penitus assignare liceat.

\InterpolationArticle{23}
Si progressio $X$ cum $\cos mx$ et $\sin mx$ non abrumpitur, valor ipsius $T$
in art. praec. inventus incompletus erit, co\"efficientesque $\gamma,\gamma',\delta'$ etc. a veris
$\alpha,\alpha',\beta'$ etc. diversi erunt. Et quidem facile perspicitur, si co\"efficientes terminorum sequentium in $X$ per eosdem characteres denotentur, quibus in art. 21
usi sumus, pro $\gamma,\gamma',\delta'$ etc. usque ad $\gamma^{m-1},\delta^{m-1}$ prorsus eosdem valores prodire,
quos illic eruimus; duos ultimos autem fieri
\[
\begin{aligned}
&\gamma^m=\frac12\alpha^m+\frac12(\alpha^m+\alpha^{3m})\cos\mu a\\[5pt]&\qquad{}+\frac12(\alpha^{3m}+\alpha^{5m})\cos2\mu a+\mathrm{etc.}\\[5pt]
&\qquad +\frac12(\beta^m+\beta^{3m})\sin\mu a+\frac12(\beta^{3m}+\beta^{5m})\sin2\mu a+\mathrm{etc.}\\[5pt]
&\delta^m=\frac12\beta^m-\frac12(\beta^m-\beta^{3m})\cos\mu a\\[5pt]&\qquad{}-\frac12(\beta^{3m}-\beta^{5m})\cos2\mu a-\mathrm{etc.}\\[5pt]
&\qquad -\frac12(\alpha^m-\alpha^{3m})\sin\mu a-\frac12(\alpha^{3m}-\alpha^{5m})\sin2\mu a-\mathrm{etc.}
\end{aligned}
\]
\SourceLeaf{310}
In hoc itaque casu co\"efficientes $\gamma,\gamma',\delta',\gamma'',\delta''\mathinner{.\,.\,.}\gamma^{m-1},\delta^{m-1}$ pro veris
$\alpha,\alpha',\beta',\alpha'',\beta''\mathinner{.\,.\,.}\alpha^{m-1},\beta^{m-1}$ adoptare licet. quatenus sequentes $\alpha^{m+1},\beta^{m+1}$ etc.
negligi possunt; ad ultimos vero $\gamma^m,\delta^m$ nisi ipsi pro negligendis haberi possunt,
haec conclusio non est extendenda.

\InterpolationArticle{24}
Per disquisitionem inde ab art. 20 traditam eo pervenimus. ut propositis $\mu$
valoribus functionis $X$ arcus indeterminati $x$, periodum completam formantibus, functionem $(X')=$
\[
\begin{aligned}
&\gamma+\gamma'\cos x+\gamma''\cos2x+\mathinner{.\,.\,.}+\gamma^m\cos mx\\[5pt]
&\quad+\delta'\sin x+\delta''\sin2x+\mathinner{.\,.\,.}+\delta^m\sin mx
\end{aligned}
\]
quae illis omnibus satisfaciat, assignare possimus, ita, ut numerus $m$, ordinem
huius functionis exprimens sit vel $=\frac12\mu-\frac12$ vel $=\frac12\mu$, prout $\mu$ impar est vel
par. Ex artt. 15, 19 concluditur, quamvis aliam similem functionem, quae iisdem valoribus satisfacit, sub forma $X'+Y\sin\frac12\mu(x-a)$ contentam esse debere,
ubi $Y$ functionem indefinitam arcus $x$ exhibet formae $G$ vel $F$ (art. 13), prout $\mu$
impar est vel par. In casu priore praeter $X'$ functio alia similis non datur, iisdem
valoribus satisfaciens, nisi ex ordine altiore, unde $X'$ cum $X$ identica erit, si
constat $X$ ordinem $m^{\mathrm{tum}}$ non egredi. In casu posteriore autem infinitae quidem
aliae similes functiones ordinis $m^{\mathrm{ti}}$ iisdem valoribus satisfaciunt, omnes tamen in
eo convenient, quod sub forma $X'+h\sin m(x-a)$ contenti erunt, designante
$h$ quantitatem definitam: quamobrem, si constat $X$ ordinem $m^{\mathrm{tum}}$ non egredi,
saltem $X$ et $X'$ aliter quam in co\"efficientibus ipsorum $\cos mx$ et $\sin mx$
non discrepabunt; in functione $X'$ autem inter hos co\"efficientes aequatio
$\gamma^m\sin ma-\delta^m\cos ma=0$ locum habebit. Porro patet, si qua functio $X''$ similis formae ut $X'$, quae iisdem valoribus satisfaciat, aliunde constet, iisdem satisfieri per functionem quamcunque $X''+Y\sin\frac12\mu(x-a)$, atque $Y$ ita determinari posse, ut haec functio ordinem $m^{\mathrm{tum}}$ non egrediatur. Ad hunc finem observamus, si qua pars ipsius $Y$ exhibeatur per $K\cos(\kappa x+k)$, designante $K$
co\"efficientem definitum, $k$ arcum definitum, $\kappa$ numerum vel integrum vel fractione
$\frac12$ affectum (prout $\mu$ par vel impar) --- ad quam formam singulae partes reduci
possunt, si cosinus et sinus eiusdem arcus contrahuntur --- pars respondens producti $Y\sin\frac12\mu(x-a)$ exhibebitur per
\SourceLeaf{311}
\[
\begin{aligned}
&\frac12K\sin((\kappa+\frac12\mu)x+k-\frac12\mu a)\\[5pt]&\qquad{}-\frac12K\sin((\kappa-\frac12\mu)x+k+\frac12\mu a)
\end{aligned}
\]
sive quod eodem redit per
\[
\begin{aligned}
&L\sin(\lambda x+l)-L\sin((\lambda-\mu)x+l+\mu a)
\end{aligned}
\]
designante $\lambda$ integrum. Quocirca in $X''$ pro quavis parte $L\sin(\lambda x+l)$ substituere licebit hanc $L\sin((\lambda-\mu)x+l+\mu a)$ et simili ratione pro hac rursus
$L\sin((\lambda-2\mu)x+l+2\mu a)$ porroque $L\sin((\lambda-3\mu)x+l+3\mu a)$ etc. Ut similis
conclusio ad cosinus extendatur. sufficit observatio, cosinus pro sinubus prodire,
si modo pro $l$ scribatur $l+90^{0}$. Hinc colligitur, si in $X''$ occurrat terminus
$L\cos\lambda x$, designante $\lambda$ integrum maiorem quam $\frac12\mu$, qui proin sub formam
$\lambda=\nu\mu\pm\lambda'$ poni potest, ita ut $\nu$ sit integer atque $\lambda'$ non maior quam $\frac12\mu$, pro
isto termino substitui posse
\[
\begin{aligned}
&L\cos(\pm\lambda'x+\nu\mu a)=L\cos\nu\mu a\cos\lambda'x\mp L\sin\nu\mu a\sin\lambda'x
\end{aligned}
\]
et similiter pro termino tali $L\sin\lambda x$ substitui poterit
\[
\begin{aligned}
&L\sin(\pm\lambda'x+\nu\mu a)=L\sin\nu\mu a\cos\lambda'x\pm L\cos\nu\mu a\sin\lambda'x
\end{aligned}
\]
Hoc modo manifesto functio $X''$ deprimetur ad aliam, ordinis certo non maioris
quam $\frac12\mu^{\mathrm{ti}}$.

Per hanc itaque operationem $X''$ semper transibit in ipsam $X'$, quoties $\mu$
impar est; in casu altero autem, ubi $\mu$ par est, certo in functionem talem, quae
tantummodo in co\"efficientibus ipsorum $\cos mx$ et $\sin mx$ a functione $X'$ diversa
esse potest. Ut ex illis co\"efficientibus co\"efficientes respondentes in $X'$ deducantur, methodus generalis art. 15 adhiberi potest, ex qua sine negotio sequitur, si
in illa functione termini ultimi sint $K\cos mx+L\sin mx$, pro his substitui debere, ut functio $X'$ prodeat,
\[
\begin{aligned}
&\quad(K\cos ma+L\sin ma)\cos ma\cos mx\\[5pt]
&+(K\cos ma+L\sin ma)\sin ma\sin mx
\end{aligned}
\]
Hi termini fiunt $=K\cos mx+L\sin mx$, si inter $K$ et $L$ aequatio $K\sin ma=L\cos ma$
locum habet, ut per calculum facile confirmatur: tunc igitur illa functio cum $X'$
omnino iam identica est.
\endgroup
